\RequirePackage[thinlines]{easybmat} %-- muss aufgrund eines Bugs vor etex und tikz geladen werden
\documentclass[a4paper, twoside, headsepline, index=totoc,toc=listof, fontsize=10pt, cleardoublepage=empty, headinclude, DIV=13, BCOR=13mm, titlepage,]{scrartcl}
\usepackage{scrtime} % KOMA, Uhrzeit ermoeglicht
\usepackage{scrpage2} % wie fancyhdr, nur optimiert auf KOMA-Skript, leicht andere Syntax
\usepackage{etoolbox}
\usepackage{letltxmacro}

% Compiler abhaengige IFs
\usepackage{ifxetex,ifluatex}
\newif\ifxetexorluatex
\ifxetex
  \xetexorluatextrue
\else
  \ifluatex
    \xetexorluatextrue
  \else
    \xetexorluatexfalse
  \fi
\fi

%--Farbdefinitionen
%-- muss vor tikz geladen werden
\usepackage[usenames, table, x11names]{xcolor}
\definecolor{dark_gray}{gray}{0.45}
\definecolor{light_gray}{gray}{0.6}

%--Zum Zeichnen
%-- muss vor polyglossia bzw. babel geladen werden
\usepackage{tikz}
\usepackage{tikz-cd}
\usetikzlibrary{external}
\tikzset{>=latex}
\usetikzlibrary{shapes,arrows,intersections}
\usetikzlibrary{calc,3d}
\usetikzlibrary{decorations.pathreplacing,decorations.markings}
\usetikzlibrary{angles}

%-- Konfiguration von tikzexternalize
\tikzexternalize[prefix=tikz/, up to date check=diff]

\ifxetexorluatex
\tikzset{external/system call={xelatex \tikzexternalcheckshellescape %-- verwende LuaLaTeX, wegen dynamischer Speicherallokation
    -halt-on-error -interaction=batchmode -jobname "\image" "\texsource"}}
\else
\tikzset{external/system call={pdflatex \tikzexternalcheckshellescape 
    -halt-on-error -interaction=batchmode -jobname "\image" "\texsource"}}
\fi

%-- tikzexternalize fuer tikzcd deaktivieren, da inkompatibel
\AtBeginEnvironment{tikzcd}{\tikzexternaldisable}
\AtEndEnvironment{tikzcd}{\tikzexternalenable}

%-- um Inkompatibilitaeten von quotes und polyglossia bzw. babel zu vermeiden
\tikzset{
  every picture/.append style={
    execute at begin picture={\shorthandoff{"}},
    execute at end picture={\shorthandon{"}}
  }
}
\usetikzlibrary{quotes}
\usepackage{pgfplots}
\usepgfplotslibrary{colormaps}
\pgfplotsset{compat=1.10}



%-- Mathe
\usepackage{mathtools} % beinhaltet amsmath
\mathtoolsset{showonlyrefs, centercolon}
\newtagform{brackets}[\textbf]{[}{]}
\usetagform{brackets}
\usepackage{wasysym}
\usepackage{amssymb} %zusätzliche Symbole
\usepackage{latexsym} %zusätzliche Symbole
\usepackage{stmaryrd} %für Blitz
\usepackage{nicefrac} %schräge Brüche
\usepackage{cancel} %Befehle zum Durchstreichen
\usepackage{mathdots}
% \usepackage{bm}
\DeclareSymbolFont{bbold}{U}{bbold}{m}{n}
\DeclareSymbolFontAlphabet{\mathbbold}{bbold}
\newcommand{\mathds}[1]{\mathbb{#1}} %Um Kompatibilitaet mit frueheren Benutzung von dsfont herzustellen
\newcommand{\ind}{\mathbbold{1}} % charakteristische-Funktion-Eins
\def\mathul#1#2{\color{#1}\underline{{\color{black}#2}}\color{black}} %farbiges Untersteichen im Mathe-Modus

%-- Alles was mit Schrift und XeteX zu tun hat
\ifxetexorluatex
    % XeLaTeX or LuaTeX
	\ifxetex
		\usepackage{mathspec} %beinhaltet fontspec 
	\else
		\usepackage[no-math]{fontspec}
	\fi
    \usepackage{polyglossia} %babel-ersatz
    \setmainlanguage[spelling=new,babelshorthands=false]{german}
	\newcommand\glqq{"}
	\newcommand\grqq{"}
	\defaultfontfeatures{Mapping=tex-text, WordSpace={1.4}} %
    \setmainfont[Ligatures=Common, BoldFont={* Bold}, ItalicFont={* Light Italic},SmallCapsFont={Linux Libertine O}, SmallCapsFeatures={Letters=SmallCaps}]{Source Sans Pro}
	\setsansfont[Scale=MatchLowercase,Ligatures=Common, BoldFont={* Medium}]{Ubuntu}
	\ifxetex
		\setallmonofonts[Scale=MatchLowercase, ItalicFont={*}]{Consolas} 
	\else
		\setmonofont[Scale=MatchLowercase, ItalicFont={*}]{Consolas}
	\fi
	\usepackage{xltxtra}
	\usepackage{fontawesome}
	\usepackage{microtype}
\else
    % default: pdfLaTeX
    \usepackage[ngerman]{babel}
    \usepackage[T1]{fontenc}
    \usepackage[utf8]{inputenc}
	\usepackage{textcomp} %verhindert ein paar Fehler bei den Fonts
    \usepackage[babel=true, tracking=true]{microtype}
	\usepackage{lmodern}
	\renewcommand{\familydefault}{\sfdefault} 
\fi


%--Mixed
\usepackage[neverdecrease]{paralist}
\usepackage[german=quotes]{csquotes}
\usepackage{makeidx}
\usepackage{booktabs}
\usepackage{wrapfig}
\usepackage{float}
\usepackage[margin=10pt, font=small, labelfont={sf, bf}, format=plain, indention=1em]{caption}
\captionsetup[wrapfigure]{name=Abb. }
\usepackage{stackrel}
\usepackage{ifthen}
\usepackage{multicol}
\usepackage{qrcode}
\flushbottom

%--Hyphenation
\hyphenation{di-men-si-o-na-le}


%--Unterstreichung
\usepackage[normalem]{ulem}
\setlength{\ULdepth}{1.8pt}

%--Indexverarbeitung
\newcommand{\bet}[1]{\textbf{#1}}
\newcommand{\Index}[1]{\textbf{#1}\index{#1}}
\makeindex
\setindexpreamble{{\noindent \itshape Die \emph{Seitenzahlen} sind mit Hyperlinks zu den entsprechenden Seiten versehen, also anklickbar} \ifxetexorluatex \faHandUp\fi \par \bigskip}
\renewcommand{\indexpagestyle}{scrheadings}


%--Marginnote & todonotes
\usepackage{marginnote}
\renewcommand*{\marginfont}{\color{dark_gray} \itshape \footnotesize }
\usepackage[textsize=small]{todonotes}

%--Todonotes mit tikz/externalize kompatibel machen
\LetLtxMacro{\oldtodo}{\todo}
\renewcommand{\todo}[2][]{\tikzexternaldisable\oldtodo[#1]{#2}\tikzexternalenable}
\LetLtxMacro{\oldmissingfigure}{\missingfigure}
\renewcommand{\missingfigure}[2][]{\tikzexternaldisable\oldmissingfigure[{#1}]{#2}\tikzexternalenable}

%--Konfiguration von Hyperref pdfstartview=FitH, 
\usepackage[hidelinks, pdfpagelabels,  bookmarksopen=true, bookmarksnumbered=true, linkcolor=black, urlcolor=SkyBlue2, plainpages=false,pagebackref, citecolor=black, hypertexnames=true, pdfauthor={Jannes Bantje}, pdfborderstyle={/S/U}, linkbordercolor=SkyBlue2, colorlinks=false]{hyperref}

% \let\orighref\href
\ifxetexorluatex
\newcommand{\hrefsym}[2]{\href{#1}{\texttt{#2\,\faExternalLink}}}
\newcommand{\hrefsymmail}[2]{\href{#1}{\texttt{\faEnvelope\,#2}}}
\renewcommand{\url}[1]{\hrefsym{#1}{\nolinkurl{#1}}}
\fi
\providecommand{\hrefsym}{\href}
% \usepackage{cleveref}
% \crefformat{footnote}{#2\footnotemark[#1]#3}

%--Römische Zahlen
\newcommand{\RM}[1]{\MakeUppercase{\romannumeral #1{}}}

%--Punkte (nach hyperref)
\usepackage{ellipsis}

%%--Abkürzungen etc.
\newcommand{\light}{\text{\Large $\lightning$}}

%-- Definitionen von weiteren Mathe-Befehlen
\DeclareMathOperator{\re}{Re} %Realteil
\DeclareMathOperator{\im}{Im} %Imaginaerteil
\DeclareMathOperator{\id}{id} %identische Abbildung
\DeclareMathOperator{\Sp}{Sp} %Spur
\DeclareMathOperator{\sgn}{sgn} %Signum
\DeclareMathOperator{\alt}{Alt} %Alternierende n-linearFormen
\DeclareMathOperator{\End}{End} %Endomorphismen
\DeclareMathOperator{\Vol}{Vol} %Volumen
\DeclareMathOperator{\dom}{dom} %Domain
\DeclareMathOperator{\Hom}{Hom} %Homomorphismen
\DeclareMathOperator{\bild}{Bild} %Bild
\DeclareMathOperator{\Kern}{Kern }
\DeclareMathOperator{\rg}{rg} %Rang
\DeclareMathOperator{\Rg}{Rg} %Rang
\DeclareMathOperator{\diam}{diam} %Durchmesser
\DeclareMathOperator{\dist}{dist} %Distanz
\DeclareMathOperator{\grad}{grad} %Gradient
\DeclareMathOperator{\dive}{div} %Gradient
\DeclareMathOperator{\rot}{rot} %Rotation
\DeclareMathOperator{\hess}{Hess} %Hesse-Matrix
\DeclareMathOperator{\koker}{Koker} %Kokern
\DeclareMathOperator{\aut}{Aut} %Automorphismen
\DeclareMathOperator{\ord}{ord} %Ordnung
\DeclareMathOperator{\ggT}{ggT} %ggT
\DeclareMathOperator{\kgV}{kgV} %kgV
\DeclareMathOperator{\Gr}{Gr} %Gerade
\DeclareMathOperator{\Kr}{Kr} %Kreis
\DeclareMathOperator{\Char}{char} %Charakteristik
\DeclareMathOperator{\Aut}{Aut} %Automorphismen
% \DeclareMathOperator{\D}{D} %formale Ableitung
\newcommand{\D}{\ensuremath{\mathrm{D}\mkern-1.5mu}}
\newcommand{\mathd}{\ensuremath{\mathrm{d}\mkern-1.0mu}}
\newcommand{\Tmap}{\ensuremath{\mathrm{T}\mkern-1.0mu}}
\newcommand{\opL}{\ensuremath{\mathrm{L}\mkern-0.6mu}}
\DeclareMathOperator{\cov}{cov} %Kovarianz
\DeclareMathOperator{\Gal}{Gal} %Galois
\DeclareMathOperator{\supp}{supp} %Träger
\DeclareMathOperator{\Sym}{Sym} %Symmetrische Gruppe
\DeclareMathOperator{\Zyl}{Zyl} %Zylinder
\DeclareMathOperator{\Mod}{Mod} %Moduln
\DeclareMathOperator{\EPK}{EPK} %Einpunktkompaktifizierung
\DeclareMathOperator{\conj}{conj} %Konjugation
\DeclareMathOperator{\ann}{ann} %Annulator
\DeclareMathOperator{\tor}{tor} %Torsionsmodul
\DeclareMathOperator{\Int}{Int} %Interior/Inneres
\DeclareMathOperator{\Br}{Br} %Brauerergruppe
\DeclareMathOperator{\GL}{GL} %allgemeine lineare Gruppe
\DeclareMathOperator{\SL}{SL} %spezielle lineare Gruppe
\DeclareMathOperator{\SU}{SU} %spezielle unitäre Gruppe
\DeclareMathOperator{\Tr}{Tr} %Spur/Trace
\DeclareMathOperator{\conv}{conv} %Konvexe Teilmenge
% \DeclareMathOperator{\tr}{tr}
\newcommand{\tr}{\mathrm{tr}}
\newcommand{\w}{\mathrm{w}}
\newcommand{\weakT}[1]{\ensuremath{\mathcal{T}_{#1}^{\mkern+1.0mu\text{\raisebox{0.4ex}{$\mathrm{w}$}}}}}
\newcommand{\weakTstar}[1]{\ensuremath{\mathcal{T}_{#1}^{\mkern+1.0mu\text{\raisebox{0.4ex}{$\mathrm{w}$}}^*}}}
\newcommand{\rand}[1]{\ensuremath{\partial^{\scriptscriptstyle #1}}}


% -- Kategorien
\DeclareMathOperator{\Ob}{Ob} %Objekte
\DeclareMathOperator{\Mor}{Mor} %Morphismen

% \DeclareMathOperator{\TOP}{TOP} %Topologische Räume
\newcommand{\TOP}{\textsc{Top}}
% \DeclareMathOperator{\HTOP}{HTOP} %Topologische Räume + Homotopieäquivalenz
\newcommand{\HTOP}{\textsc{HTop}}
% \DeclareMathOperator{\VR}{VR} %Vektorräume
\newcommand{\VR}{\textsc{VR}}
% \DeclareMathOperator{\MOD}{MOD} %Moduln
\newcommand{\MOD}{\textsc{Mod}}
% \DeclareMathOperator{\MENGEN}{MENGEN} %Mengen
\newcommand{\MENGEN}{\textsc{Mengen}}
% \DeclareMathOperator{\MAN}{MAN} %Mannigfaltigkeiten
\newcommand{\MAN}{\textsc{Man}}
% \DeclareMathOperator{\GRUPPEN}{GRUPPEN} %Gruppen
\newcommand{\GRUPPEN}{\textsc{Gruppen}}
% \DeclareMathOperator{\ABELGRUPPEN}{ABEL.GRUPPEN} %abelsche Gruppen
\newcommand{\ABELGRUPPEN}{\textsc{Abel.Gruppen}}
% \DeclareMathOperator{\KAT}{KAT} %Kategorien
\newcommand{\KAT}{\textsc{Kat}}
% \DeclareMathOperator{\FUN}{FUN} %Funktoren
\newcommand{\FUN}{\textsc{Fun}}


\newcommand{\op}{\mathrm{op}}

\newcommand{\Ker}{\ensuremath{\Kern \,}}
\newcommand{\zirkel}{\rotatebox[origin = c]{-90}{$\varangle$}}
\newcommand{\wirkung}{\!\curvearrowright\!}

\newcommand{\oE}{\OE}

%--Beweisende
\newcommand{\bewende}{\ifmmode \tag*{$\square$} \else \hfill $\square$ \fi}

%--Volumen
\newcommand{\vol}[1]{\ensuremath{\Vol_{#1}}}

%--Skalarprodukt
\DeclarePairedDelimiterX\sprod[2]{\langle}{\rangle}{#1\,\delimsize\vert\,#2}

%--Betrag, Gaußklammer
\DeclarePairedDelimiter{\abs}{\lvert}{\rvert}
\DeclarePairedDelimiter{\ceil}{\lceil}{\rceil}

%--Norm
\DeclarePairedDelimiter\doppelstrich{\Vert}{\Vert}
\newcommand{\norm}[2][\relax]{
\ifx#1\relax \ensuremath{\doppelstrich*{#2}}
\else \ensuremath{\doppelstrich*{#2}_{#1}}
\fi}


%--Umklammern
\DeclarePairedDelimiter\enbrace{(}{)}
\DeclarePairedDelimiter\benbrace{[}{]}
\DeclarePairedDelimiter\lenbrace{<}{>}

%--Mengen
\DeclarePairedDelimiterX\mengenA[1]{\lbrace}{\rbrace}{#1}
\DeclarePairedDelimiterX\mengenB[2]{\lbrace}{\rbrace}{#1\, \delimsize\vert \, #2}

\makeatletter
\newcommand{\set}[2][\relax]{
\ifx#1\relax \ensuremath{
\mengenA*{#2}}
\else \ensuremath{%
  \mengenB*{#1}{#2}}
\fi}
\makeatother


%--offen und abgeschlossen
\newcommand{\off}{\! \stackrel[\text{offen}]{}{\subset} \!}
\newcommand{\abg}{\! \stackrel[\text{abg.}]{}{\subset} \!}

%--Differential
\newcommand{\diff}[2]{\ensuremath{\frac{{\partial #1}}{{\partial #2}} }}
\newcommand{\diffd}[2]{\ensuremath{\frac{\mathd #1}{\mathd #2} }}

%--Diagonale Linien für Matrizen
\newcommand\tikzmark[1]{%
  \tikzexternaldisable\tikz[overlay,remember picture,baseline] \node [anchor=base] (#1) {};\tikzexternalenable}

\newcommand\MyLine[3][]{%
  \tikzexternaldisable\begin{tikzpicture}[overlay,remember picture]
    \draw[#1] (#2.south) -- (#3.north);
  \end{tikzpicture}\tikzexternalenable}
  
%--Jordankasten
\newcommand{\Jord}{ \tikzexternaldisable \begin{BMAT}{|ccc|}{|ccc|} %
				0 & & \\ %
				1\tikzmark{z} &  \diagdown & \\ %
				& \tikzmark{u}1 \MyLine[thick]{z}{u} &  0 %
			\end{BMAT} \tikzexternalenable}
			
%--Jordankasten mit Lambda
\newcommand{\JordLa}[1][\empty]{%
	\tikzexternaldisable\ifthenelse{\equal{#1}{\empty}}
		{\begin{BMAT}{|ccc|}{|ccc|} %
				\lambda & & \\ %
				1\tikzmark{z} &  \diagdown & \\ %
				& \tikzmark{u}1 \MyLine[thick]{z}{u} &  \lambda  %
			\end{BMAT}}
		{\begin{BMAT}{|ccc|}{|ccc|} %
				\lambda_{#1}  & & \\ %
				1\tikzmark{z} &  \diagdown & \\ %
				& \tikzmark{u}1 \MyLine[thick]{z}{u} &  \lambda_{#1}  %
			\end{BMAT}}\tikzexternalenable}

%--Inhaltsverzeichnis
\usepackage[tocindentauto]{tocstyle}
\usetocstyle{KOMAlike}			

\newcommand{\fach}{Differentialgeometrie \RM{1}.}
\newcommand{\semester}{WiSe 2014}
\newcommand{\homepage}{https://wwwmath.uni-muenster.de/u/frederik.witt/diffgeoI.html}
\newcommand{\prof}{Frederik Witt}
\newcommand{\verfasser}{Arne Grauer}
\providecommand{\verfasser}{Jannes Bantje}

%--Konfiguration von scrheadings
\setheadsepline{1pt}[\color{light_gray}]
\pagestyle{scrheadings}
%\fancyhf[H, F]{}
\clearscrheadfoot

\providecommand{\shortFach}{\fach}
\lehead{ \includegraphics[height=0.5 cm,keepaspectratio]{../!config/Bilder/Logo_WWU_Muenster_light_gray.pdf}}
\rehead{\rule{0cm}{0.5cm}\footnotesize \sffamily \color{light_gray} \verfasser -- Mitschrift \shortFach}
\lohead{\rule{0cm}{0.5cm} \footnotesize \sffamily \color{light_gray} Stand: \today \; \thistime[:]}
\rohead{\includegraphics[height=0.5 cm,keepaspectratio]{../!config/Bilder/fb10logo_gray.pdf}}


\ofoot[{ \color{dark_gray} \LARGE \sffamily \thepage}]{{ \color{dark_gray} \LARGE \sffamily \thepage}} %hier wir auch der plain Stil bearbeitet!
\automark{section}
\ifoot{ \color{dark_gray} \small \leftmark}

%--Metadaten
\author{\verfasser}
\titlehead{\includegraphics[height=1.5cm, keepaspectratio]{../!config/Bilder/Logo_WWU_Muenster.pdf}%
\hfill \includegraphics[height=1.3cm, keepaspectratio]{../!config/Bilder/fb10logo.pdf}}
\title{Skript \fach}
\subtitle{Mitschrift der Vorlesung  \enquote{\fach} von \prof}

\publishers{
Aktuelle Version verfügbar bei: \bigskip\\
\normalsize
\begin{minipage}[t]{0.45\textwidth}
	\centering{
	\qrcode[height=3.5cm, version=4]{https://github.com/JaMeZ-B/latex-wwu}
	\medskip\\
	\includegraphics[height=0.4cm, keepaspectratio]{../!config/Bilder/github_octo.pdf} 
	\includegraphics[height=0.4cm, keepaspectratio]{../!config/Bilder/GitHub_Logo.pdf} (inklusive Sourcecode)\\
	{\small\url{https://github.com/JaMeZ-B/latex-wwu}}}
\end{minipage}
\hfill
\begin{minipage}[t]{0.45\textwidth}
	\centering{
	\qrcode[height=3.5cm, version=4]{btsync://B6WH2DISQ5QVYIRYIEZSF4ZR2IDVKPN3I?n=latex_share}
	\medskip\\
	\includegraphics[height=0.4cm, keepaspectratio]{../!config/Bilder/bt_sync_logo.pdf} 
	{\large \textbf{Bittorrent} Sync} \\
	\texttt{B6WH2DISQ5QVYIRYIEZSF4ZR2IDVKPN3I}
	}
\end{minipage}
}



\begin{document}
\pagenumbering{Roman}
\maketitle
\begin{abstract}
\section*{Vorwort --- Mitarbeit am Skript}
Dieses Dokument ist eine Mitschrift aus der Vorlesung \enquote{\fach, \semester}, gelesen von \prof. Der Inhalt entspricht weitestgehend dem Tafelanschrieb. Für die
Korrektheit des Inhalts übernehme ich keinerlei Garantie! Für Bemerkungen und Korrekturen -- und seien es nur Rechtschreibfehler -- bin ich sehr dankbar. 
Korrekturen lassen sich prinzipiell auf drei Wegen einreichen: 
\begin{itemize}
	\item Persönliches Ansprechen in der Uni, Mails an \hrefsymmail{mailto:\mail}{\mail} (gerne auch mit annotieren PDFs) 
	\item \emph{Direktes} Mitarbeiten am Skript: Den Quellcode poste ich auf GitHub (siehe oben), also stehen vielfältige Möglichkeiten der Zusammenarbeit zur Verfügung:
	Zum Beispiel durch Kommentare am Code über die Website und die Kombination Fork + Pull Request. Wer sich verdient macht oder ein Skript zu einer Vorlesung, die 
	ich nicht besuche, beisteuern will, dem gewähre ich gerne auch Schreibzugriff.
	
	Beachten sollte man dabei, dass dazu ein Account bei \url{github.com} notwendig ist, der allerdings ohne Angabe von persönlichen Daten angelegt werden kann. 
	Wer bei GitHub (bzw. dem zugrunde liegenden Open-Source-Programm \enquote{\texttt{git}}) -- verständlicherweise -- Hilfe beim Einstieg braucht, dem helfe ich gerne 
	weiter. Es gibt aber auch zahlreiche empfehlenswerte Tutorials im Internet.\footnote{zB. \url{https://try.github.io/levels/1/challenges/1}, ist auf Englisch, aber dafür 
	interaktives LearningByDoing}
	\item \emph{Indirektes} Mitarbeiten: \TeX-Dateien per Mail verschicken. 
	
	Dies ist nur dann sinnvoll, wenn man einen ganzen Abschnitt ändern möchte (zB. einen alternativen Beweis geben), da ich die Änderungen dann per Hand einbauen muss! Ich freue mich aber auch über solche Beiträge!
\end{itemize}
\section*{Vorlesungshomepage}
{\centering 
\begin{minipage}[c][][c]{\textwidth}
	\centering \qrcode[height=3cm]{\homepage} \medskip\\
	\footnotesize \url{\homepage}
\end{minipage}
\par}

\section*{Über die Differentialgeometrie}
\begin{description}
	\item[1. Differentialgeometrie] Mannigfaltigkeiten $\Rightarrow $ Differentialtopologie
	\item[2. Riemannsche Geometrie] Längen und Winkel messen $\leadsto$ Geodäte (lokal) kürzeste Verbindung zwischen 2 Punkten einer Mannigfaltigkeit.
	
	Krümmung
	
	Allgemeine Relativitätstheorie: Krümmung = Gravitation ist für solche Effekte verantwortlich
\end{description}
\end{abstract}

\setcounter{page}{1}
\tableofcontents
\cleardoubleoddemptypage
\pagenumbering{arabic}
\setcounter{page}{1}


\section{Differenzierbare Mannigfaltigkeiten} % (fold)
\label{sec:1}
Konvention: glatt = $C^\infty$.\\
Wiederholung: Diffeomorphismus $f : U \subseteq \mathds{R}^n \to \mathds{R}^n$ glatt, bijektiv und $f ^{-1}$ glatt.

\subsection{Untermannigfaltigkeiten} % (fold)
\label{sub:11}

\subsubsection[Definition: Submersion, Immersion und Étale]{Definition} % (fold)
\label{ssub:111}
Eine Abbildung $f : U \subseteq \mathds{R}^n \to \mathds{R}^m$ glatt heißt \Index{Submersion}/\Index{Immersion}/\Index{Étale} in $x \in U$, falls
$\mathd_x f : \mathds{R}^n \to \mathds{R}^m$ surjektiv/injektiv/Isomorphismus ist. \marginnote{Insbesondere $n\ge m$/$n \le m$/$n=m$}. $f$ heißt schlechthin
\bet{Submersion}/\bet{Immersion}/\bet{Étale}, falls $f$ Submersion/Immersion/Étale für alle $x \in U$ ist.
% subsubsection 111 (end)

\subsubsection[Bemerkung: Submersiv/immersiv/étale sind offene Bedingungen]{Bemerkung} % (fold)
\label{ssub:112}
Submersiv/immersiv/étale sind offene Bedingungen, d.h. ist $f$ Submersion/Immersion/Étale in $x$, so auch in einer hinreichend kleinen Umgebung von $x$ 
(vergleiche auch Satz 15 unten)
% subsubsection 112 (end)

\subsubsection{Satz über inverse Funktionen (inverse function theorem)} % (fold)
\label{ssub:113}
Sei $f : U \subseteq \mathds{R}^n \to \mathds{R}^m$ étale in $x$. Dann existiert eine offene Menge $x \in V \subset U \subset \mathds{R}^n$, sodass $f\big|_{V}$ ein
Diffeomorphismus auf seinem Bild ist, d.h. $f\big|_{V} : V \xrightarrow{\sim} f(V) $ ist bijektiv und $f\big|_V$, $(f\big|_V)^{-1}$ glatt.
% subsubsection 113 (end)

\subsubsection{Satz über implizite Funktionen (implicit function theorem)} % (fold)
\label{ssub:114}
Sei $f : U \subset \mathds{R}^{n+k} \to \mathds{R}^n$ eine Submersion in $0$ mit $f(0)=0$. Dann existiert eine glatte Funktion $g$ auf einer Umgebung von $0 \in \mathds{R}^n$
mit Werten in $\mathds{R}^k$, d.h.
\[
	g=g(x_1, \ldots , x_n) = \begin{pmatrix}
		g_1(x_1, \ldots ,x_n) \\
		\vdots \\
		g_k(x_1, \ldots ,x_n)
	\end{pmatrix}
\]
mit
\[
	f \enbrace[\big]{x_1, \ldots ,x_n , g_1(x_1, \ldots , x_n), \ldots , g_k(x_1, \ldots ,x_n)} = f \enbrace[\big]{x,g(x)} = 0
\]
% subsubsection 114 (end)

\subsubsection{Satz (Normalenformen von Submersionen/Immersionen)} % (fold)
\label{ssub:115}
\begin{enumerate}[a)]
	\item Sei $f : U \subseteq \mathds{R}^{n+k} \to \mathds{R}^n$ eine Submersion mit $f(0)=0$. Dann existiert ein Diffeomorphismus 
	$\varphi : 0 \in V \subset \mathds{R}^{n+k} \to \mathds{R}^{n+k}$ mit $\varphi(V) \subset U$ und 
	$f \circ \varphi(x_1, \ldots , x_{n+k}) = (x_1, \ldots , x_n) = \pi_{\mathds{R}^n}(x_1, \ldots , x_{n+k})$, d.h. modulo einer Diffeomorphismus ist jede Submersion 
	äquivalent zu einer Projektion.
	\item Sei $f : U \subset \mathds{R}^n \to \mathds{R}^{n+k}$ eine Immersion in $0$, $f(0)=0$. Dann existiert ein Diffeomorphismus 
	\[
		\varphi : 0 \in V \subset \mathds{R}^n \to \mathds{R}^{n+k}
	\]
	mit $\varphi(0)=0$ und 
	\[
		\varphi \circ f (x_1, \ldots ,x _n) = \enbrace*{x_1, \ldots , x_n, 0 , \ldots , 0} = i(x_1, \ldots , x_n) 
	\]
	d.h. modulo Diffeomorphismen ist $f$ äquivalent zur kanonischen Einbettung von $\mathds{R}^n \hookrightarrow \mathds{R}^{n+k}$.
\end{enumerate}
\minisec{Beweis}
\begin{enumerate}[(i)]
	\item Wähle Koordinaten $x_1, \ldots , x_n$ mit $\mathd_0 f = (A \, B)$ mit $A \in \mathds{R}^{n \times n}$, $B \in \mathds{R}^{k \times n}$ mit $\det A \not= 0$.
	Definiere $F : U \subset \mathds{R}^n \times \mathds{R}^k \to \mathds{R}^n \times \mathds{R}^k $ durch $(p,q) \mapsto \enbrace*{f(p,q), q}$. Dann gilt
	\[
		\mathd_{(0,0)} F = \begin{pmatrix}
			\mathd_0 f & \\
			0 & \id_{k \times k}
		\end{pmatrix} = \begin{pmatrix}
			A & B \\
			0 & \id_{k \times k} 
		\end{pmatrix}
		\Longrightarrow \det \mathd_0 F \not= 0
	\]
	Mit \ref{ssub:113} folgt: $\varphi = F ^{-1}$ Diffeomorphismus um $0 \in \mathds{R}^{n+k}$ herum.
	\item Definiere $F : U \times \mathds{R}^k \to \mathds{R}^n \times \mathds{R}^k$. $F(p,q) = f(p)+ (0,q)$.
	\[
		\mathd_{(0,0)} F = \begin{pmatrix}
			\mathd_0 f & 0 \\
			0 & \id_{k \times k}
		\end{pmatrix},
	\]
	also $\det \mathd_{(0,0)} F \not= 0$. $\Rightarrow  \exists \varphi = F ^{-1}$ um $0$ mit $\varphi(f(p))=\varphi \circ F (p,0) = (p,0) = i(p)$ \bewende
\end{enumerate}
\minisec{Bemerkung}
Durch Komposition mit Translationen können Submersionen/Immersionen mit $F(x)=y$ (nicht notwendigerweise $=0$) behandelt werden.
% subsubsection 115 (end)

\subsubsection[Definition: Untermannigfaltigkeit]{Definition} % (fold)
\label{ssub:116}
$M \subset \mathds{R}^{n+k}$ heißt $n$-dimensionale \Index{Untermannigfaltigkeit}, falls $\forall x \in M : \exists x \in U \subset \mathds{R}^{n+k}$ Umgebung, sowie eine
Submersion $f : U \to \mathds{R}^k$ mit $U \cap M = f ^{-1}(0)$. \marginnote{Verallgemeinerung des $n$-dimensionalen Unterraumes}
% subsubsection 116 (end)

\subsubsection[Beispiele für Untermannigfaltigkeiten]{Beispiele} % (fold)
\label{ssub:117}
\begin{enumerate}[(i)]
	\item Die $n$-\Index{Sphäre} $S^n := \set[x=(x_1, \ldots ,x_{n+1}) \in \mathds{R}^{n+1}]{\abs*{x}^2=1 } $ ist eine $n$-dimensionale Untermannigfaltigkeit des 
	$\mathds{R}^{n+1}$. Definiere $f : \mathds{R}^{n+1} \to \mathds{R}$ durch $f(x_1, \ldots , x_{n+1}) = x_1^2 + \ldots + x_{n+1}^2 -1$. $\Rightarrow f ^{-1}(0) = S^n$ nach
	Definition.
	\[
		\mathd_x f = (2 x_1, \ldots , 2 x_{n+1}) : \mathds{R}^{n+1} \to \mathds{R}
	\]
	\item Die hyperbolischen Räume $H_c^n := \set[x=(x_1, \ldots ,x_{n+1}) \in \mathds{R}^{n+1}]{-x_1^2 + \sum_{i=2}^{n+1}x_i^2 = c}$\begin{figure}[H]
\centering{
\begin{tikzpicture}[xscale=0.6, yscale=0.8]
	    \begin{axis}[view={40}{20}]
	       \addplot3[surf,
	       colormap/cool,
	       samples=60,
	       domain=-2*pi:2*pi,y domain=-2*pi:2*pi,
	       z buffer=sort,
		   shader=faceted,]
			({0.5*sqrt(x^2+1)*cos(deg(y))},
			 {0.5*sqrt(x^2+1)*sin(deg(y))},
			{x});
	    \end{axis}
\end{tikzpicture}
\begin{tikzpicture}[xscale=0.6, yscale=0.8]
	    \begin{axis}[view={40}{20}]
	       \addplot3[surf,
	       colormap/cool,
	       samples=60,
	       domain=1:2*pi,y domain=-2*pi:2*pi,
	       z buffer=sort,
		   shader=faceted,]
			({0.5*sqrt(x^2-1)*cos(deg(y))},
			 {0.5*sqrt(x^2-1)*sin(deg(y))},
			{x});
			\addplot3[surf,
	       colormap/cool,
	       samples=60,
	       domain=-2*pi:-1,y domain=-2*pi:2*pi,
	       z buffer=sort,
		   shader=faceted,]
			({0.5*sqrt(x^2-1)*cos(deg(y))},
			 {0.5*sqrt(x^2-1)*sin(deg(y))},
			{x});
	    \end{axis}
\end{tikzpicture}
\begin{tikzpicture}[xscale=0.6, yscale=0.8]
	    \begin{axis}[view={40}{20}]
	       \addplot3[surf,
	       colormap/cool,
	       samples=60,
	       domain=-2*pi:2*pi,y domain=-2*pi:2*pi,
	       z buffer=sort,
		   shader=faceted,]
			({0.5*sqrt(x^2)*cos(deg(y))},
			 {0.5*sqrt(x^2)*sin(deg(y))},
			{x});
	    \end{axis}
\end{tikzpicture}
\caption{hyperbolischer Raum für $c>0$, $c<0$, $c=0$}}
\end{figure}
	Sei $f:\mathds{R}^{n+1} \to \mathds{R}$ definiert durch $f_c(x) = -x_1^2 + \sum_{i=2}^{n+1}x_i^2 - c$.\\
	$\Rightarrow$ $f$ ist glatt und $f_c^{-1}(0) = H_c^n$\\
	Es bleibt zu zeigen, das $f$ eine Submersion ist.
	\[
		\mathd_x f = (-2x_1,2x_2, \ldots,2x_{n+1}): \mathds{R}^{n+1} \to \mathds{R}	
	\]
	Es gilt $\mathd_x f(\nu) = 0 \forall \nu \in \mathds{R}^{n+1} \Leftrightarrow x_i = 0\ \forall i=1, \ldots,n+1$.\\
Das heißt $\exists i: x_i \neq 0 \Rightarrow \mathd_x f\ \text{Submersion}$, wobei $\exists i: x_i \neq 0$ gilt, falls $x \in H_c^n, c \neq 0$. Also ist für $c \neq 0$ $H_c^n$ eine $n$-dimensionale Untermannigfaltigkeit. Für $c = 0$ ist $H_0^n \backslash \set{0}$ eine $n$-dimensionale Untermannigfaltigkeit.
	\item Der $n$-Torus $T^n = \set[x=(x_1,x_2, \ldots,x_{2n-1},x_{2n})\in \mathds{R}^{2n}]{x_1^2 + x_2^2 = 1, \ldots, x_{2n-1}^2 + x_{2n}^2 = 1}$
	\[
		f:\mathds{R}^{2n} \to \mathds{R},\ f(x) = 
		\begin{pmatrix}
			x_1^2 + x_2^2 - 1\\
			\vdots\\
			x_{2n-1}^2 + x_{2n}^2 - 1		
		\end{pmatrix}
	\]
	$\Rightarrow$ $f$ glatt, $f^{-1}(0)=T^n$
	\[
		\mathd_x f =
		\begin{pmatrix}
			2x_1 & 2x_2 & 0 & 0 &\ldots & 0\\
			0 & 0 & 2x_3 &2x_4 &  &0\\
			\vdots &&& \ddots &\ddots & \\
			0 &&&& 2x_{2n-1} & 2x_{2n}
		\end{pmatrix}
	\]
	$\Rightarrow$ $\mathd_x f$ ist eine Submersion.
	\item Die orthogonale Gruppe $O(n) \subset \mathds{R}^{n \times n}$ ist eine Untermannigfaltigkeit der Dimension $\frac{n(n-1)}{2}$.
	\[
		O(n) = \set[A \in \mathds{R}^{n \times n}]{AA^T = \id_{\mathds{R}^n} = f^{-1}(0)},\ f(A) = AA^T - \id_{\mathds{R}^n}
	\]
	Es gilt $f(A)^T = {A^T}^T \cdot A^T - \id^T = AA^T - \id = f(A)$. Also $f:\mathds{R}^{n \times n} \to \text{Sym}(\mathds{R}^{n \times n}) = \set[A \subseteq \mathds{R}^{n \times n}]{A^T = A}$.\\
	$\Rightarrow$ $f$ glatt, $f^{-1}(0) = O(n)$
	\[
		\mathd_A f:\mathds{R}^{n \times n} \to \text{Sym}(\mathds{R}^{n \times n}),\ \mathd_A f = HA^T+AH^T
	\]
	Sei $S \in \text{Sym}(\mathds{R}^{n \times})$. Setze $H = \frac{SA}{2} \Rightarrow \mathd_A f(\frac{SA}{2}) = \frac{1}{2}(SAA^T+ASA^T) = \frac{1}{2}(S +  S) = S$\\
	$\Rightarrow$ $O(n)$ ist eine Untermannigfaltigkeit
\end{enumerate}
% subsubsection 117 (end)

\subsubsection[Bemerkung: ]{Bemerkung} % (fold)
\label{ssub:118}
$M$ kann von der Form $M = f^{-1}(0)$ sein, $f$ glatt, aber keine Submersion. Zum Beispiel
	\[
		M = \set[(x,x)\in \mathds{R}^2]{x \in \mathds{R}} = f^{-1}(0),\ f(x,y) = x^3 +y^3	
	\]
%subsubsection118 (end)
	
\subsubsection[Satz: lokale Parametrisierung von Untermannigfaltigkeiten]{Satz}
\label{ssub:119}
Es sind äquivalent:
\begin{enumerate}[(i)]
	\item $M^n\subset \mathds{R}^{n+k}$ ist $n$-dim Untermannigfaltigkeit
	\item $\forall x \in M^n\ \exists$ Umgebungen $U$ und $V \subset \mathds{R}^{n+k}$ von $x$ und $0$, und ein Diffeomorphismus $\phi: U \to V$ mit $\phi(M\cap U) \to \underbrace{(\mathds{R}^n \times \set{0}\cap V)}_{=: W}$
	\item $\forall x \in M^n\ \exists$ Umgebungen $U\subset \mathds{R}^{n+k}$ und $W\subset \mathds{R}^n$, sowie eine glatte Abbildung $g:W \to \mathds{R}^{n+k}$, sodass $g$ ein Homöomorphismus von $W$ auf $M\cap U$ und eine Immersion ist. $g$ nennt man eine \Index{lokale Parametrisierung}
\end{enumerate}
\missingfigure{Beweis des Satzes durch Zeichnung}
%subsubsection119 (end)

\subsubsection[Beispiele zur lokalen Parametrisierung]{Beispiele}
\label{ssub:1110}
\begin{enumerate}[(i)]
	\item Der $2$-Torus $T^2$ ist das Bild von
	\[
		g:\mathds{R}^2 \to \mathds{R}^2,\ g(x,y) = \enbrace[\big]{\cos(x),\sin(x),\cos(y),\sin(y)}\subset T^2 \subset \mathds{R}^4
	\]
	$g$ ist lokale Parametrisierung:
	\[
		\mathd_{(x,y)} g = 
		\begin{pmatrix}
		-\sin(x) & 0\\
		\cos(x) & 0\\
		0 & -\sin(x)\\
		0 & \cos(y)
		\end{pmatrix}
	\]
	$\mathd_{(x,y)} g$ hat somit immer vollen Rang 2, also ist $\mathd_{(x,y)} g$ injektiv. Es folgt somit, dass $g$ eine Immersion ist.
%	\begin{wrapfigure}{R}{7cm}
%		\begin{minipage}{6cm}
%			Sei $X$ ein topologischer Raum , $\sim$ eine Äquivalenzrelation. So ist $\nicefrac{X}{\sim}$ ein topologischer Raum versehen mit der Quotiententopologie, also $U\subset\ \nicefrac{X}{\sim}$ offen $\Leftrightarrow\ p^{-1}(U)$ offen. Somit folgt für $X \overset{p}{\to} \nicefrac{X}{\sim}$, dass $p$ stetig und auch ein lokaler Hoöomorphismus ist.
%		\end{minipage}
%	\end{wrapfigure}
	\missingfigure{Diagramm zum lok. Homöomorphismus von g}
	$\hat{g}$ ist stetig (Quotiententopologie), $\hat{g}$ bijektiv und $T^2$ kompakt. Damit folgt, dass $\hat{g}$ ein Homöomorphimus und somit $g$ ein lokaler Homöomorphismus ist.
	\item 
	\[
		g:\mathds{R}^2 \to \mathds{R}^3\text{, } g(x,y) = \enbrace[\big]{\sin(x)\cos(y),\sin(x)\sin(y),\cos(x)} \in S^2\subset \mathds{R}^3
	\]
	$g$ ist ein lokaler Homöomorphismus, jedoch ist $\mathd_{(x,y)}$ nicht injektiv, das heißt keine lokale Parametrisierung nahe der Pole. Also müssen wir eine andere Parametrisierung wählen, zum Beispiel $g(x,y) = \enbrace{x,y,\sqrt{1-x^2-y^2}}$
\end{enumerate}
%subsubsection1110 (end)

\subsubsection[Bemerkung]{Bemerkung}
\label{ssub:1111}
\begin{enumerate}[(i)]
	\item 
		Betrachte:
		\begin{figure}[H]
		\centering{
			\begin{tikzpicture}[scale = 0.8]
				%\draw[help lines] (0,0) grid (12,4);
				\draw (1,2) node[below]{$\infty$} -- (3,2) node[below]{$\mathds{R}$} -- (5,2) node[below]{$-\infty$};
				\draw[color = blue, thick, ->] (2,2) -- (3,2) node[above]{$\nu,\ \mathd_x g \neq 0$} -- (4,2);
				\draw[->] (5.5,2) .. controls +(25:0.25) and +(180:0.25) .. (6,2.25) node[above]{$g$} .. controls +(0:0.25) and +(155:0.25) .. (6.5,2);
				\draw[<->] (7,4) -- (7,0) -- (11,0);
				\draw[<->] (9,2) .. controls +(90:1) and +(180:0.5) .. (10,3.5)
				.. controls +(0:0.5) and +(90:0.5) .. (10.75,2.75)
				.. controls +(270:0.5) and +(0:1) .. (9.5,2)
				.. controls +(180:0.5) and +(0:0.5) .. (8.5,2)
				.. controls +(180:1) and +(90:0.5) .. (7.25,1.25)
				.. controls +(270:0.5) and +(180:0.5) .. (8,0.5)
				.. controls +(0:0.5) and +(270:1) .. (9,2);
				\draw[color = blue, thick, ->] (9.5,2) .. controls +(0:1) and +(270:0.5) .. (10.75,2.75) node[right]{$\nu$};
			\end{tikzpicture}
			}		
		\end{figure}
		$g$ ist immersiv, hat aber keine Umkehrfunktion und ist somit kein Homöomophismus.
		\item
		\missingfigure{Weiteres Beispiel zur lok. Par.}
		$g$ ist keine lokale Parametrisierung, da $g$ nicht immersiv in $0$ ist. $g(t) = \enbrace[\big]{g_1(t),g_2(t)}\in \mathds{R}^2 \Rightarrow g_1^2(t) = g_2^3(t) \Rightarrow \mathd_0 g = 0$
\end{enumerate}
%subsubsection1111 (end)

\subsubsection[Satz: Wechsel lok. Parametrisierungen]{Satz}
\label{ssub:1112}
Sei $M^n$ eine $n$-dim glatte Untermannigfaltigkeit, $x\in M$, und seien $U_1,U_2$ Umgebungen von $x$ in $\mathds{R}^{k+n}$ mit lok. Parametrisierung $g_1$ und $g_2$ von $U_1\cap M$ und $U_2\cap M$ definiert auf $\Omega_1$ und $\Omega_2$ (das heißt $g_i: \Omega_i \to \mathds{R}^{n+k}$). Dann ist
\[
	g_2^{-1} \circ g_1:\Omega_1 \cap g_1^{-1}(U_2) \overset{g_1}{\to} g_1(\Omega_1) \cap U_2 \overset{g_2^{-1}}{\to} \Omega_2 \cap g_2^{-1}(U_1)
\]
ein $C^{\infty}$-Diffeomorphismus.
\missingfigure{Veranschaulichung des Wechsel lok. Parametrisierung}
\minisec{Bemerkung}
	\begin{itemize}
		\item $h:\Omega_1 \to \Omega_2$ Homöomorphismus, offene Menge in $\Omega_1$ $\leftrightarrow$ offene Menge in $\Omega_2$
		\item $\varphi:\Omega_1 \to \Omega_2$ Diffeomorphismus, falls Homöom. und $C^{\infty}(\Omega_1)$ $\leftrightarrow$ $C^{\infty}(\Omega_2)$
	\end{itemize}
%subsubsection1112 (end)
% subsection 11 (end)

\subsection{Abstrakte Mannigfaltigkeit}
\label{sub:12}

\subsubsection[Definition: Karte und Atlas]{Definition}
\label{ssub:121}
Sei $M$ ein topologischer Raum. Ein \Index{Atlas} ist eine Familie $\set{(U_i,\varphi_i)}_i\in I$ mit
\begin{enumerate}[(i)]
	\item $U_i \subset M$ offen, $M = \bigcup_{i \in I} U_i$ ($\set{U_i}$ offene Überdeckung)
	\item $\set{\varphi_i:U_i \to \Omega_i \subset \mathds{R}^n}$ ist eine Familie von Homöomorphismen mit:
	\[
		U_{ij} = U_i \cap U_j \neq \emptyset \text{, dann ist } \varphi_{ij} = \varphi_i \circ \varphi_j^{-1}: \varphi_j(U_{ij}) \to \varphi_i(U_{ij})
	\]
	ein $C^{\infty}$-Diffeomorphismus.
\end{enumerate}
Ein Paar $(U_i,\varphi_i)$ heißt \Index{Karte} und $\varphi_{ij}$ heißt Übergangsfunktion oder \index{Kartenwechsel}.
%subsubsection 121(end)

\subsubsection[Definition: Abstrakte differenzierbare Mannigfaltigkeit]{Definiton}
\label{ssub:122}
Sei $M$ ein topologischer Raum.
\begin{enumerate}[(i)]
	\item Zwei Atlanten $\mathcal{A}_1 = \set{(U_i,\varphi_i)}_{i\in I}$ und $\mathcal{A}_2 = \set{(V_j,\psi_j)}_{j\in J}$ sind äquivalent, genau dann wenn $\forall i\in I,j\in J$ mit $U_i\cap V_j\neq \emptyset$
	\[
		\varphi_i\circ \psi_j^{-1}:\psi(U_i\cap V_j) \to \varphi(U_i\cap V_j)
	\]
	ein $C^{\infty}$-Diffeomorphismus, das heißt die Vereinigung und $\mathcal{A}_1 \cup \mathcal{A}_2$ ist wieder ein Atlas von $M$.
	\item \Index{Diffenrenzierbare Struktur} auf $M$ ist die Äquivalenzklasse eines Atlanten. Ein Hausdorffraum $M$ mit einer differenzierbaren Struktur heißt \Index{(abstrakte) differenzierbare Mannigfaltigkeit} oder Mannigfaltigkeit.
	\item Nehmen alle Karten eines Atlanten Werte in $\mathds{R}^n$ an, so heißt $n$ die Dimension der Mannigfaltigkeit.
\end{enumerate}
% subsubsection 122(end)

\subsubsection[Beispiele für differenzierbare Mannigfaltigkeiten]{Beispiele}	
\label{ssub:123}
\begin{enumerate}[(i)]
	\item $\mathds{R}^n$ ist Mannigfaltigkeit der Dimension $n$, Atlas: $\set{\mathds{R}^n,\id_{\mathds{R}^n}}$ und $U\subset \mathds{R}^n$ ist Mannigfaltigkeit der Dimension $n$, Atlas: $\set{U,\id_{\mathds{R}^n}\big\vert_U}$
	\item Untermannigfaltigkeit im Sinne von Abschnitt $1.1$ sind abstrakte Mannigfaltigkeiten.
	\item Mannigfaltigkeit $N\subset M$ ist eine Untermannigfaltigkeit von $M$, falls für alle $y\in N$ eine Karte $(U,\varphi)$ von M um $y$ existiert, sodass $\varphi(U\cap N)$ Untermannigfaltigkeit von $\mathds{R}^n$ im Sinne von Abschnitt $1.1$ ist.\\
Dann gilt zum Beispiel $U \overset{offen}{\subset} \mathds{R}^n$ ist Untermannigfaltigkeit im "abstrakten" Sinne. Untermannigfaltigkeit $N\subset \mathds{R}^n$ im Sinne von $1.1$ ist Untermannigfaltigkeit im abstrakten Sinne.
	\item Die $n$-Sphäre, z.B $n = 1$
	\begin{figure}[H]
		\tikzsetnextfilename{stereographic_projection}
		\centering{\begin{tikzpicture}[scale=0.8]
			\draw[name path=real] (-5.2,0) node[above right]{$\mathds{R}^n$} -- (5.2,0);
			\draw[name path=sphere, thick] (0,0) circle[radius=2cm];
			
			%einzelner Punkt mit Beschriftung
			\draw[name path=g1,lightgray, thick, dashed] (0,2) -- (-4,0);
			\path[name intersections={of=g1 and real, by=in1}];
			\path[name intersections={of=g1 and sphere}];
			\draw[fill=SteelBlue2] (in1) circle[radius=0.06cm] node[below, color=SteelBlue3]{$h^+(p)$};
			\draw[fill=SteelBlue2] (intersection-2) circle[radius=0.06cm] node[above left, color=SteelBlue3]{$p$};
			
			\begin{scope}
			\foreach \x in {-3, -2.2, 2, 2.5, 3.4, 5} %
			{	
				\draw[name path=g1,lightgray, thick, dashed] (0,2) -- (\x,0);
				\path[name intersections={of=g1 and real, by=in1}];
				\path[name intersections={of=g1 and sphere}];
				\draw[fill=SteelBlue2] (in1) circle[radius=0.06cm];
				\draw[fill=SteelBlue2] (intersection-1) circle[radius=0.06cm];
				\draw[fill=SteelBlue2] (intersection-2) circle[radius=0.06cm];
			}
			\end{scope}
			\foreach \x in {-3.1,-2.1, -0.7, 0.2, 0.7, 1.8, 2.9} %
			{	
				\path[name path=g1] (0,2) -- (\x,-2);
				\path[name intersections={of=g1 and real, by=in1}];
				\path[name intersections={of=g1 and sphere}];
				\draw[name path=g1,lightgray, thick, dashed] (0,2) -- (intersection-2);
				\draw[fill=SteelBlue2] (in1) circle[radius=0.06cm];
				\draw[fill=SteelBlue2] (intersection-2) circle[radius=0.06cm];
			}
			\draw[fill=OrangeRed2] (0,2) circle[radius=0.07cm] node[above]{$N$};
			\draw[fill=OrangeRed2] (0,-2) circle[radius=0.07cm] node[below]{$S$};
		\end{tikzpicture}
		\caption{Stereografische Projektion}}
	\end{figure}
		Allgemein: $U_N = S^n\backslash \set{N},\ U_S = S^n\backslash \set{S}\subset \mathds{R}^{n+1}$
		\begin{align*}
			\varphi_N:\ U_N \to \mathds{R}^n,\ \enbrace{x_1, \ldots , x_{n+1}} \mapsto \frac{1}{1-x_{n+1}} \enbrace{x_1, \ldots, x_n}\\
			\varphi_S:\ U_S \to \mathds{R}^n,\ \enbrace{x_1, \ldots , x_{n+1}} \mapsto \frac{1}{1+x_{n+1}} \enbrace{x_1, \ldots, x_n}\\
			\varphi_S \circ \varphi_N^{-1}:\ \mathds{R}^n\backslash \set{0} \to \mathds{R}^n\backslash \set{0},\ x\in \mathds{R}^n\backslash \set{0} \mapsto \frac{x}{\|x\|^2}\\
		\end{align*}
		\minisec{Bemerkung} $\set{\enbrace*{S^n,\id_{\mathds{R}^{n+1}}\vert_{S^n}}}$ ist kein Atlas, denn $S^n$ ist nicht offen im $\mathds{R}^{n+1}$.
		\item Der $n$-Torus $T^n$. Betrachte $f:\mathds{R}^n \to \mathds{C}^n,\ \enbrace{x_1, \ldots, x_n} \mapsto \enbrace{\exp(ix_1), \ldots, \exp(ix_n)}\in \underbrace{S^1\times \ldots \times S^1}_{n-\text{mal}} \subset \mathds{C}\times \ldots \times \mathds{C}$, so folgt $f(\mathds{R}^n) = T^n$
		\missingfigure{Veranschaulichung der Wahl der Karten}
		Schränke $f\vert_{\enbrace{a_1-\epsilon,a_1+\epsilon}\times \ldots \times \enbrace{a_n-\epsilon,a_n+\epsilon}}$ auf sein Bild ein. $\varphi_a := f^{-1} : \underset{= U_a \subset T}{f(I_a)} \to I_a \overset{\text{offen}}{\subset} \mathds{R}^n$. Somit ist $\set{\enbrace{U_a,\varphi_a}}_{a\in \mathds{R}^n}$ eine Familie von Karten, die $T^n$ überdeckt. Es bleibt zu zeigen, dass $\varphi_{ab} = \varphi_a \circ \varphi_b$ ein Diffeomorphismus ist, was wir an dieser Stelle dem Leser überlassen. Somit folgt, dass $\set{\enbrace{U_a,\varphi_a}}_{a\in \mathds{R}^n}$ ein Atlas ist.
		\item Betrachte den Hyperbolischen Raum $H^n$.
		\missingfigure{Veranschaulichung der Wahl der Karten}
		$H^n$ ist eine Mannigfaltigkeit mit einer Karte, gegeben durch stereographische Projektion. Man sagt daher, dass $H^n$ diffeomorph zu $\mathds{R}^n$
		\item Würfel $Q\subset \mathds{R}^n$
		\begin{figure}[H]
			\centering{			
			\begin{tikzpicture}[scale = 0.8]
				%\draw[help lines] (0,0) grid (6,6);
				\draw[<->] (0,3) -- (6,3);
				\draw[<->] (3,0) -- (3,6);
				\draw (1,1) -- (5,1) -- (5,5) -- (1,5) -- (1,1);
				\draw (3,3) node[below right]{$S^n$} circle(2);
				\foreach \x/\y in {1/2,1/4,2/5,4/5,5/2,5/4,4/1,2/1}
				{
				\draw[dashed] (\x,\y) -- (3,3);
				}
				\node at (2,3.6) {$\sigma$};
			\end{tikzpicture}	
			\caption{Der Würfel $Q$ ist keine Untermannigfaltigkeit des $\mathds{R}^n$!}
			}
		\end{figure}
		$Q$ ist jedoch eine abstrakte Mannigfaltigkeit, denn $Q$ ist ein Hausdorffraum und bijektiv zu $S^n$ durch Projektion $sigma$. So ist $\mathcal{A} := \set{\enbrace*{\sigma^{-1}(U_N),\varphi_N \circ \sigma},\enbrace*{\sigma^{-1}(U_S),\varphi_S \circ \sigma}}$ ein Atlas, da $\varphi_2 \circ \varphi_1^{-1} = \varphi_S \circ \sigma \circ \sigma^{-1} \circ \varphi_N$ und somit ist $Q$ eine abstrakte Mannigfaltigkeit.
		\item Sei $\mathds{K} = \mathds{R},\ \mathds{C}$. So ist der \Index{projektive Raum}
		\begin{align*}
			\mathds{K}\mathds{P}^n =& \text{Menge der Ursprungsgraden in $\mathds{K}^{n+1}$}\\
			=& \nicefrac{\mathds{K}^{n+1}\backslash \set{0}}{\sim},\ x,y\in \mathds{K}^{n+1}\backslash \set{0}\ \text{seien äquivalent},\ x\sim y \Leftrightarrow \exists \lambda\in \mathds{K}\backslash \set{0}: x = \lambda y\\
			=& \nicefrac{\mathds{K}^{n+1}\backslash \set{0}}{\mathds{K}^*}\ \text{wobei}\\
			&[x]\in \nicefrac{\mathds{K}^{n+1}\backslash \set{0}}{\mathds{K}^*} \mapsto\ <x>_{\mathds{K}}\ \text{lineare Erzeugnis von $x$ in $\mathds{K}^{n+1}$}
		\end{align*}
		Wir schreiben $[x] = [x_0:\ldots:x_n]$, falls $x=(x_0,\ldots,x_n)\in \mathds{K}^{n+1}$, "homogene Koordinaten".
		\minisec{Bemerkung} $[x_0:\ldots:x_n] = [\lambda x_0:\ldots:\lambda x_n],\ \lambda \in \mathds{K}^{n+1}$
		\begin{enumerate}[(1)]
			\item Betrachte $p:\ x\in \mathds{K}^{n+1}\backslash \set{0} \mapsto [x] = p(x)\in \nicefrac{\mathds{K}^{n+1}\backslash \set{0}}{\sim}$ so folgt, dass $\mathds{K}\mathds{P}^n$ versehen mit der Quotiententopologie ein Hausdorff-Raum ist.
			\begin{figure}[H]
			\centering{
	\begin{tikzpicture}[scale=1]
		\draw[->] (0,0) -- (2,0) node[right]{$x_2$};
		\draw (0,0) -- (0,1);
		\draw[->] (0,0) -- (0,0,2) node[left]{$x_1$};
		\draw (0.25,-0.5,-0.5) -- (-0.5,1,1);
		\draw (-0.15,-0.5,-0.3) -- (0.3,1,0.6);
		\draw (-0.3,-0.5,0.35) -- (0.6,1,-0.7);
		\draw (0.55,-0.5,0.15) -- (-1.1,1,-0.3);
		%\draw ( ,-0.5, ) -- ( ,1, );
		\draw[SeaGreen4, thick, fill=SeaGreen4!20, fill opacity=0.4] (2,1,2) -- (2,1,-2) -- (-2,1,-2) -- (-2,1,2) -- cycle;
		\draw (-0.5,1,1) -- (-1.25,2.5,2.5);
		\draw (0.3,1,0.6) -- (0.75,2.5,1.5);
		\draw (0.6,1,-0.7) -- (1.5,2.5,-1.75);
		\draw (-1.1,1,-0.3) -- (-2.75,2.5,-0.75);
		\draw[fill, Firebrick2] (-0.5,1,1) circle[radius=0.03];
		\draw[fill, Firebrick2] (0.3,1,0.6) circle[radius=0.03];
		\draw[fill, Firebrick2] (0.6,1,-0.7) circle[radius=0.03];
		\draw[fill, Firebrick2] (-1.1,1,-0.3) circle[radius=0.03];
		\draw[->] (0,1) -- (0,2) node[left]{$x_3$};
		\draw[fill, SeaGreen4] (0,1) circle[radius=0.03];
		\node[below right,SeaGreen4] at (2,1,0) {$\mathds{R}^2\times\set{1}$} ;
		%\draw (0,1) -- (0,2);
	\end{tikzpicture}
	\caption{Der 2-dimensionale projektive Raum $\mathds{R}P^2$ (projektive Ebene)}
	\begin{tikzpicture}[scale=0.8]
		\draw[->, thick] (-1,0) -- (3,0) node[right]{$x_1$};
		\draw[->, thick] (0,-0.2) -- (0,1.5) node[above]{$x_2$};
		\draw[name path=ebene, thick, SeaGreen4] (-1.5,1) -- (3,1) node[below]{$\mathds{R}\times \set{1}$};
		\foreach \phi / \rad in {30 / 3, 55 / 2, 78 / 1.7, 109 / 1.7, 134 / 2}
		{
			\draw[name path=gerade] (\phi+180:0.3) -- (0,0) -- (\phi:\rad);
			\path[name intersections={of=ebene and gerade, by=in1}];
			\draw[fill=Firebrick2] (in1) circle[radius=0.05cm];
		}	
	\end{tikzpicture}
	\caption{Der eindimensionale projektive Raum $\mathds{R}P^1$}}
	\end{figure}
	Aus Abbildung 3 bzw. 2 folgt somit $\mathds{R}\mathds{P}^n = \mathds{R} \cup \set{\infty}  \cong S^2$ und analog $\mathds{R}^2 \cup \set{N} \cong S^2$. Um zu zeigen, dass $\mathds{K}\mathds{P}^n$ eine abstrakte Mannigfaltigkeit ist, muss die Existenz eines Atlanten gezeigt werden. Betrachte $U_i = \set[[x_0:\ldots:x_n]\in \mathds{K}\mathds{P}^n]{x_i \neq 0}$. So folgt, dass $\bigcup_{i=0}^n = \mathds{K}\mathds{P}^n$, da $(0,\ldots,0)\not \in \mathds{K}^{n+1}\backslash \set{0}$, also $\not \exists: [0:\ldots:0]$. Da $[x_0:\ldots:x_n]\in U_i \Rightarrow x_i \neq 0$, ist $\varphi_i: [x_0:\ldots:x_n]=[\frac{x_0}{x_i}:\ldots:\frac{x_n}{x_i}] \overset{\varphi_i}{\to} [\frac{x_0}{x_i}:\ldots:\hat{\frac{x_i}{x_i}}:\ldots\frac{x_n}{x_i}]$ wohldefiniert.
	Es bleibt als Übung zu zeigen, dass $\varphi_i \circ \varphi^{-1}_j$ ein Diffeomorphismus ist.
		\end{enumerate}
\end{enumerate}	
% subsubsection 123(end)

\subsubsection[Bemerkungen zu abstrakten Mannigfaltigkeiten]{Bemerkung}
\label{ssub:124}
\begin{enumerate}[(i)]
	\item Die Vereinigung aller äquivalenten Atlanten ist wieder ein Atlas, der sogenannte maximale Atlas (eindeutig bestimmt). In der Praxis arbeitet man mit konkreten Atlanten.
	\item Allgemein kann man $C^p$-Atlanten betrachten, das heißt $\varphi_{ij}$ $C^p$-Diffeomorphismen $\to$ $C^p$-Mannigfaltigkeiten. Ebenso kann man Karten mit allgemeinen Zielräumen betrachten, wo der Begriff der Differenzierbarkeit/Diffeomorphismus Sinn hat. Zum Beispiel Banach- und Hilbertraum $\to$ Banach- und Hilbertmannigfaltigkeiten oder Frécheraum $\to$ Fréchemannigfaltigkeiten.
	\item Sei $M$ ein topologischer Raum mit differenzierbarer Struktur, so ist er nicht automatisch Hausdorff! Zum Beispiel
\begin{figure}[H]
\centering{
			\begin{tikzpicture}
				%\draw[help lines] (0,0) grid (8,5);
				\draw (1,4) -- (7,4);
				\draw[thick] (4,3.9) -- (4,4.1) node[below right]{$0$};
				\filldraw[SteelBlue4] (4,4.5) circle[radius=0.04cm];
				\filldraw[OrangeRed2] (4,3.5) circle[radius=0.04cm];
				
				\draw[color = SteelBlue4] (1,1.5) -- (7,1.5) node[below right]{$\mathds{R}$};
				\draw[color = OrangeRed2] (1,1) -- (7,1) node[below right]{$\mathds{R}$};
				\draw[thick,->] (4,2) -- (4,2.5) node[right]{$\nicefrac{\mathds{R}\cup\mathds{R}}{\sim}$ mit $(x,0)\sim (0,y) \Leftrightarrow x = y \neq 0$} -- (4,3);
			\end{tikzpicture}
		\caption{Beispiel einer Nichthaudorffmenge}}
\end{figure}
Benutze Atlas $\set{\enbrace{U_1=\set[(x,0)]{x\in \mathds{R}},\id},\enbrace{U_2=\set[(0,y)]{y\in \mathds{R}},\id}}$. Dies ist ein Atlas, da $U_1\cap U_2 = M\backslash \set{:}$.
\item Wir nehmen künftig zwei zusätzliche Bedingungen für Mannigfaltigkeiten an:
	\begin{enumerate}[(a)]
		\item $M$ zusammenhängend $\to$ jede Mannigfaltigkeit hat wohldefinierte Dimension schreibe $M^n$
		\item $M$ erfüllt das 2. Abzählbarkeitsaxiom. Für uns: $\exists$ abzählbarer Atlas.
	\end{enumerate}
\end{enumerate}
%subsubsection 124(end)

% subsection 12(end)

\subsection{Glatte Abbildung}
\label{sub:13}

\subsubsection[Definition: Glatte Funktionen zwischen zwei Mannigfaltigkeiten]{Definition}
\label{ssub:131}

Seien $M^n$, $N^m$ zwei Mannigfaltigkeiten und sei $f:M^n\to N^m$ eine stetige Abbildung. Dann heißt $f$ glatt, falls für alle $x$ Karten $\enbrace{U,\varphi} \subset M$ und $\enbrace{V,\psi} \subset N$ um $x$ beziehungsweise $f^{-1}(x)$ existieren mit $f(U)\subset V$ und 
\[
	\psi \circ f \circ \varphi^{-1}:\ \underset{\subset \mathds{R}^n}{\varphi(U)} \to \underset{\subset \mathds{R}^n}{\psi(V)}
\]
glatt ist. ("$f$ in Karten gelesen ist glatt")
\missingfigure{Veranschaulichung der glatten Funktionen}
Wir schreiben $C^{\infty}(M,N) = \set[f:M\to N]{f\ \text{glatt}}$ und als Spezialfall $C^{\infty}(M) = C^{\infty}(M,\mathds{R})$.
% subsubsection 131(end)

\subsubsection{Bemerkung}
\label{ssub:132}
Glattheit ist unabhängig von der Wahl der Karten! ($\leadsto$ Diffeomorphismusbedingung!)
% subsubsection 132(end)

\subsubsection{Definition}
\label{ssub:133}
\begin{enumerate}[(i)]
\item Eine glatte Abbildung $f:M\to N$ heißt Submersion/Immersion/étale, falls für ein Paar (und damit alle Karten) $\psi \circ f \circ \varphi^{-1}$ eine Submersion/Immersion/étale ist.
\item $f$ ist Diffeomorphismus zwischen $M$ und $N$ ($M$ und $N$ sind "diffeomorph"), falls $f$ bijektiv ist und $f$ sowie $f^{-1}$ glatt ist.
\item $f:M\to N$ heißt Einbettung, falls $f$ Immersion und Homöomorphismus auf sein Bild ist.
\end{enumerate}
% subsubsection 133(end)

\subsubsection{Beispiele}
\label{ssub:134}
\begin{enumerate}[(i)]
\item Wenn $M$ eine Untermannigfaltigkeit des $\mathds{R}^{n+k}$, $f:\mathds{R}^{n+k}\to \mathds{R}^m$ glatte Funktion, dann ist auch $f\vert_M$ glatt.
\begin{itemize}
\item $\set{\mathds{R}^n,\id}$, $\set{\mathds{R}^m,\id}$ $\Rightarrow\ f:\mathds{R}^n\to \mathds{R}^m$ glatt im Sinne von Definition $\ref{ssub:131}$ $\Leftrightarrow$ glatt im üblichen Sinne
\item $x\in M$, Karte um $x$: $\varphi: U\subset \mathds{R}^{n+k} \underset{\text{Diffeo.}}{\to} V = \varphi(U)\subset \mathds{R}^{n+k}$ mit $\varphi(U\cap M) = U\cap (\mathds{R}^n\times \set{0})$ (aus $\ref{ssub:119}$), so gilt
\[
f\circ \varphi^{-1}: V\cap (\mathds{R}^n\times \set{0}) \xrightarrow{\varphi^{-1} \vert_{V\cap (\mathds{R}^n\times \set{0})}\in C^{\infty}(\mathds{R}^n,\mathds{R}^{n+k})} 
M\cap U \xrightarrow{f \in C^{\infty}(\mathds{R}^{n+k},\mathds{R}^m)} \mathds{R}^m
\]
Somit folgt, dass $f\vert_M$ glatt ist. 
\missingfigure{Veranschaulichung von f}
\end{itemize}
\item $S^2$ ist diffeomorph zu $\mathds{C}\mathds{P}^1 = \set{\text{Ursprungsgeraden in}\ \mathds{C}^2} = \nicefrac{\mathds{C}\backslash \set{0}}{\mathds{C}^*}$. Definiere 
\[
f: S^2\subset \mathds{R}^3 \to \mathds{C}\mathds{P}^1,\ \enbrace{x,y,z} \mapsto
\begin{cases}
z\neq -1:\ [\frac{x+iy}{1+z}:1]\\
z = -1:\ [1:0]
\end{cases}
\]
Somit folgt
\begin{enumerate}[(1)]
\item $f$ ist bijektiv
\item Wir "lesen nun in Karten"(, wobei $U_1 = \set[ z_1:z_2 ]{z_2\neq 0} \overset{\psi_2}{\to} \frac{z_1}{z_2}$:
\[
\psi \circ f \circ \varphi_S^{-1}(x,y) = \underset{\in U_1}{\psi_2 [x+iy:1]} = x + iy \underset{\mathds{C} \simeq \mathds{R}}{=} (x,y)
\]
Also ist $\psi \circ f \circ \varphi_S^{-1} = \id$ und somit $f \vert_{S^2\backslash\set{S}}$ ist $C^{\infty}$ und $f^{-1}$ auch $C^{\infty}$. Nun brauchen wie nur noch eine Karte um S!. Zum Beispiel $\varphi_N$. Hier $\varphi_N(S) = (0,0)$. Es gilt:
\[
	f\circ \varphi_N^{-1}(x,y) = 
	\begin{cases}
		(x,y) \neq (0,0): [\varphi_S \circ \varphi_N^{-1}(x,y):1 \overset{\mathds{C} \simeq \mathds{R}}{=} [\frac{x+iy}{x^2+y^2}:1] \in U_0]\\
	[1:0] \in U_0 
	\end{cases}
\]
Somit folgt:
\[
	\psi_1 \circ f \circ \varphi_N^{-1}(x,y) =
	\begin{cases}
		\frac{x^2+y^2}{x+iy} = x-iy\\
		0
	\end{cases}
	= (x,-y)\ \text{glatt}
\]
\end{enumerate}
\end{enumerate}
% subsubsection 134(end)

\subsubsection{Definition}
\label{ssub:135}
\begin{enumerate}[(i)]
	\item Eine Mannigfaltigkeit erfüllt das zweite Abzählbarkeitsaxiom, falls es einen abzählbaren Atlas besitzt.
	\item Sein $M^n$ eine Mannigfaltigkeit. Eine Familie $\set{f_i:M\to \mathds{R}}$ glatter Funktionen heißt \Index{Zerlegung der Eins}, falls
	\begin{itemize}
		\item $\supp f_i = \overline{\set[x\in M]{f_i(x) \neq 0}}$ "Träger von $f_i$" sind kompakt und bilden "lokal endliche Familie" (das heißt, sei $K\subset M$ kompakt, so gilt $\supp f_i \cap K = \emptyset$ außer für endlich viele $i\in I$)
		\item $\forall x \sum_{i\in I} f_i(x) = 1$ (wohldefiniert, also $\sum$ endlich, da lokal endlich)
	\end{itemize}
\end{enumerate}
% subsubsection 135(end)

\subsubsection{Theorem}
\label{ssub:136}
Sei $\set{U_i}_{i\in I}$ eine Überdeckung von $M$. Ist $M$ zweifach abzählbar, so existiert eine Zerlegung der Eins, die $\mathcal{U}$ untergeordnet ist, das heißt $\supp f_i \subset U_i$. Dies ist nützlich um beliebig definierte Objekte zu globalen Objekten zu verkleben.
\missingfigure{Untergeordnete Zerlegung der Eins}
% subsubsection 136(end)

\subsubsection{Whitney Theorem}
\label{ssub:137}
Sei $M^n$ eine zweifach abzählbare Mannigfaltigkeit. Dann existiert eine Einbettung $f:M^n\to \mathds{R}^{2n}$ von M (bzw. Immersion $f:M^n\to \mathds{R}^{2n-1}$) als Untermannigfaltigkeit von $\mathds{R}^{2n}$ (bzw. $\mathds{R}^{2n-1}$)
% subsubsection 137(end)

\subsubsection{Bemerkung}
\label{ssub:138}
\begin{enumerate}[(i)]
\item Können jede "abstrakte" Mannigfaltigkit $M^n$ als eine Untermannigfaltigkeit des $\mathds{R}^{2n-1}$ realisieren.
\minisec{Aber}
Die Einbettung ist nicht kanonisch (oder "natürlich"), das heißt gewisse Eigenschaften der Untermannigfaltigkeit entstehen durch die Einbettung und sind keine intrinsische Einbettung von $M$!
\item Die Dimension $2n-1$ des Bildraumes ist optimal im folgenden Sinne
\[
\exists\ \text{Mannigfaltigkeit}\ M^n,\ \text{die sich nicht in}\ \mathds{R}^N\ \text{mit}\ N<2n-1\ \text{einbetten lassen.}
\]
\end{enumerate}
\minisec{Beweisskizze}
Wir nehmen an, das die Mannigfaltigkeit $M$ kompakt ist. So folgt mit \ref{ssub:135}, das ein endlicher Atlas $\set{\enbrace{U_i,\varphi_i : U_i \to \Omega_i\subset \mathds{R}^n}}_{i=1}^n$ mit untergeordneter Zerlegung der Eins $\set{f_i:M\to \mathds{R}}_{i=1}^n,\ \supp f_i \subset U_i$ existiert. Definiere $\psi_i:M\to \mathds{R}^n$ durch $\psi_i(x) = \varphi_i(x)f_i(x)$, so folgt $\psi_i(x) \equiv 0$ für $x\in U_i^c$ und $\psi_i$ ist glatt. Dann definiert $\psi: M\to \mathds{R}^{nr+r},\ \psi(x)=\enbrace{\psi_1(x),\ldots,\psi_r(x),f_1(x),\ldots,f_r(x)}$ eine Einbettung.\\
Von jetzt an ist $M^n$ stets eine glatte zweifach abzählbare Mannigfaltigkeit.
% subsubsection 138(end)

% subsection 14(end)

\subsection{Das Tangentialbündel einer Untermannigfaltigkeit}
\label{sub:14}

\subsubsection{Definition}
\label{ssub:141}
Sei $M^n\subset \mathds{R}^{n+k}$ eine $n$-dimensionale Untermannigfaltigkeit. Sei $x\in M$. $v\in \mathds{R}^{n+k}$ ist tangential an $M$ in $x$, falls eine $C^{\infty}$-Abbildung $c: 0\in I\subset \mathds{R} \to \mathds{R}^{n+k}$ existiert mit
\begin{itemize}
\item $c(I)\subset M$
\item $c(0) = x$
\item $c'(t) = v$
\end{itemize}
\minisec{Notation}
$T_x M^n = $ Menge der Tangentialvektoren $v$ in $x$.
% subsubsection 141(end)

\subsubsection[Satz: Vektorraumstruktur des Tangentialraumes]{Satz}
\label{ssub:142}
$T_x M^n$ ist ein $n$-dim Vektorraum$\vert_{\mathds{R}}$ ($\cong \mathds{R}^n$)
\minisec{Beweis}
Sei $M$ eine Untermannigfaltigkeit, $x\in M$, so existiert Umgebung $U\subset \mathds{R}^{n+k}$ offen, $\varphi:U\to V\subset \mathds{R}^{n+k}$ diffeomorph mit $\varphi(U\cap M) = V\cap(\mathds{R}^n\times \set{0})$. Sei $c: I\to \mathds{R}^{n+k}$ ein Weg mit $c(I)\subset M\cap U$, so ist $\varphi \circ c$ ein Weg im $V\cap(\mathds{R}^n\times \set{0})$. Es folgt, dass $T_x M \cong \mathd_{\varphi(x)} \varphi^{-1}(\mathds{R}^n\times \set{0})$, das heißt $T_x M$ übernimmt die Vektorraumstruktur von $\mathds{R}^n \times \set{0}$ 
\missingfigure{Beispiel eines Tangentialraumes am $S^2$}
% subsubsection 142

\subsubsection{Satz}
\label{ssub:142}
Sei $M^n\subset \mathds{R}^{n+k}$ eine Untermannigfaltigkeit der Dimension $n$, $x\in M$ und $U\subset \mathds{R}^{n+k}$ offene Umgebung um $x$.
\begin{enumerate}[(i)]
\item $U\cap M = f^{-1}(0)$ für $f:\mathds{R}^{n+k} \to \mathds{R}^k$ eine Submersion, so ist $T_x M^n = \ker \mathd_x f$.\\
Sei zum Beispiel $M = S^2$. So ist $f(x,y,z)=x^2+y^2+z^2-1$ eine Submersion mit $S^2 = f^{-1}(0)$. Betrachte nun den Tangentialraum am Nordpol $N$. Es gilt $\mathd_N f = (0,0,z): \mathds{R}^3 \to \mathds{R}$, also gilt für den Tangentialraum: 
\[
T_N S^2 = \ker \mathd_N f = \set{
\begin{pmatrix}
x\\
y\\
0
\end{pmatrix}
\in \mathds{R}^3}
\]
\item Sei $\varphi: U\to V\subset \mathds{R}^{n+k}$ ein Diffeomorphismus mit $\varphi(U\cap M) = V\cap (\mathds{R}^n \times \set{0})$, dann ist $T_x M = \mathd_{\varphi(x)} \varphi^{-1}(\mathds{R}^n \times \set{0})$
\item Ist $\enbrace{\Omega,g}$ eine Parametrisierung von $U\cap M$, so ist der Tangentialraum $T_{g(y)} M = \mathd_y g(\mathds{R}^n)$
\end{enumerate} 
\minisec{Beweis}
\begin{enumerate}[(i)]
\item Sei $v\in T_x M$, das heißt $v = c^(0),\ c(0) = x$. So gilt $f\circ c(t) = 0$, da $c(t)\subset M\cap U = f^{-1}(0)$. Somit folgt mit der Kettenregel, das $\mathd_{c(0)} f \circ c'(0) = \mathd_x f(v) = 0$. Also ist $v\in \ker d_x f$ und es gilt $T_x M \subset \ker d_x f$. Da $T_x M$ und $\ker \mathd_x f$ $n$-dimensional, folgt somit die Gleichheit.
\item Siehe \ref{ssub:142}
\item Sei $\tilde{c}\subset \Omega$ eine Kurve, $c := g\circ \tilde{c}$ Kurve durch $c(0) = g\enbrace*{\tilde{c}(0)}\in M$. So ist $c'(0) = d_{\tilde{c}(0)} g(\tilde{c'(0)})\in T_{c(0)} M$. Also gilt $d_y g(\mathds{R}^n) \subset T_{g(y)} M$ und die Gleichheit folgt wie in (i) durch die Dimension.
\end{enumerate}
% subsubsection 142(end)
% subsection 14(end)


\subsection{Das Tangentialbündel einer abstrakten Mannigfaltigkeit}
\label{sub:15}

\subsubsection{Definition}
\label{ssub:151}
Sei $M^n$ eine abstrakte Mannigfaltigkeit, $x\in M$. Ein \Index{Tangentialvektor} von $M$ in $x$ ist eine Äquivalenzklasse von Kurven $\alpha : I\to M$ mit $\alpha(0) = x$. Dabei heißen $\alpha:I\to M$ und $\beta:J\to M$ äquivalent, falls eine Karte $(U,\varphi)$ um $x$ existiert mit
\[
\enbrace{\varphi \circ \alpha}'(0) = \enbrace{\varphi \circ \beta}'(0)
\]
\missingfigure{Veranschaulichung der Definition}
\minisec{Notation}
$T_x M =$ Raum der Tangentialvektoren
% subsubsection 151(end)

\subsubsection{Bemerkung}
\label{ssub:152}
\begin{enumerate}[(i)]
\item Gilt die oben genannte Äquivalenzbedingung für eine Karte, dann gilt sie für alle.
\minisec{Beweis}
Sei $\enbrace{V,\psi}$ eine andere Karte. Es gilt
\begin{align*}
\enbrace{\psi \circ \alpha}'(0) &= \enbrace{\psi \circ \varphi^{-1} \varphi \circ \alpha}'(0)\\
&= \mathd_{\varphi(\alpha(0))} (\psi \circ \varphi^{-1})\enbrace[\big]{(\varphi \circ \alpha)'(0)}\\
&= \mathd_{\varphi(\beta(0))} (\psi \circ \varphi^{-1})\enbrace[\big]{(\varphi \circ \beta)'(0)} = \enbrace{\psi \circ \beta}'(0)
\end{align*}
\item $T_x M$ ist $n$-dimensionaler Vektorraum$\vert_{\mathds{R}}$.
\minisec{Beweis}
Sei $\enbrace{U,\varphi}$ eine Karte um $x$. Sei $\Gamma_{(U,\varphi,x)}: T_x M\to \mathds{R}^n$, $\underset{\in T_x M}{[\alpha]} \mapsto (\varphi \circ \alpha)' (0)$. Aus der Äquivalenzrelation folgt, das $\Gamma$ wohldefiniert und bijektiv. Also können wir die Vektorraumstruktur auf $T_x M$ durch $\Gamma_{(U,\varphi,x)}$ erklären. Dies ist wohldefiniert, da $\mathd_{\varphi(x)} (\psi \circ \varphi^{-1})$ für jede weitere Karte $(V,\psi)$ um $x$ ein linearer Isomorphismus ist. (Da $\psi \circ \varphi^{-1}$ Diffeomorphismus.) 
\end{enumerate}
% subsubsection 152(end)

\subsubsection{Definiiton}
\label{ssub:153}
Sei $M$ eine Mannigfaltigkeit, $x\in M$, $(U,\varphi)$ und $(V,\psi)$ zwei Karten um $x$. Seien $u,v\in \mathds{R}^n$. Als Äquivalenzrelation ist definiert:
\[
u \sim v \xLeftrightarrow{\text{Def.}} v = \mathd_{\varphi(x)}(\psi \circ \varphi^{-1})(u)
\]
Ein \Index{Tangentialvektor} in $x$ auf $M$ ist eine solche Äquivalenzklasse $[(U,\varphi,u)]$
\missingfigure{}
% subsubsection 153(end)

\subsubsection{Bemerkung}
\label{ssub:154}
Sei $x\in M$, $(U,\varphi)$ Karte um $x$. So ist
\[
\Theta_{(U,\varphi,x)}: \underset{\in \varphi(x) + \mathds{R}^n}{u} \to [(U,\varphi,u)]\in T_x M
\]
eine Bijektion, die wohldefiniert $\text{VR}\vert_{\mathds{R}^n}$ auf $T_x M$ induziert.
% subsubsection 154(end)

\subsubsection{Definition}
\label{ssub:155}
$TM := \bigcap_{x\in M} T_x M$ ist das \Index{Tangentialbündel} von $M$.

% section 1 (end)



\cleardoubleoddemptypage
\pagenumbering{Alph}
\setcounter{page}{1}
\printindex
\listoffigures
\end{document}
